% CV of Jungsuk Oh. Build with script/build-cv (XeLaTeX + Inter).
\documentclass[10pt]{article}

\usepackage[letterpaper, top=0.55in, bottom=0.5in, left=0.65in, right=0.65in]{geometry}
\usepackage{fontspec}
\usepackage{xcolor}
\usepackage{titlesec}
\usepackage{enumitem}
\usepackage{tabularx}
\usepackage[hidelinks]{hyperref}

\setmainfont{Inter}[
  Extension      = .otf,
  UprightFont    = *-Regular,
  ItalicFont     = *-Italic,
  BoldFont       = *-SemiBold,
  BoldItalicFont = *-SemiBoldItalic,
  Numbers        = Lining,
]
\newfontfamily\displayfont{InterDisplay}[Extension = .otf, UprightFont = *-SemiBold]

\definecolor{accent}{HTML}{1D4ED8}
\definecolor{ink}{HTML}{0F172A}
\definecolor{muted}{HTML}{475569}
\definecolor{rule}{HTML}{CBD5E1}

\hypersetup{
  pdftitle   = {Jungsuk Oh - CV},
  pdfauthor  = {Jungsuk Oh},
  pdfsubject = {Curriculum Vitae},
}

\pagestyle{empty}
\setlength{\parindent}{0pt}
\raggedright
\hyphenpenalty=10000
\exhyphenpenalty=10000
\fontsize{9.6}{12.4}\selectfont
\color{ink}

% Section heading: small accent caps with a thin rule.
\titleformat{\section}
  {\color{accent}\footnotesize\bfseries\addfontfeature{LetterSpace=10}}
  {}{0pt}{\MakeUppercase}[\vspace{1pt}{\color{rule}\titlerule}]
\titlespacing*{\section}{0pt}{9pt}{5pt}

\setlist[itemize]{leftmargin=1.1em, itemsep=0.5pt, topsep=1.5pt, parsep=0pt, label={\color{muted}\textbullet}}
\newlist{pubs}{enumerate}{1}
\setlist[pubs]{leftmargin=2.5em, labelsep=0.6em, align=left, itemsep=5pt, topsep=2pt, parsep=0pt, font=\color{muted}}

\newcommand{\link}[2]{\href{#1}{\color{accent}#2}}
\newcommand{\when}[1]{\hfill{\color{muted}#1}}
\newcommand{\sub}[1]{{\color{muted}#1}}

\begin{document}
\fontsize{9.6}{12.4}\selectfont

% ---------------------------------------------------------------- header
{\displayfont\fontsize{23}{27}\selectfont Jungsuk Oh}\par
\vspace{2pt}
\sub{Ph.D. Student, Graduate School of Data Science, Seoul National University}\par
\vspace{2pt}
\link{mailto:luke0112@snu.ac.kr}{luke0112@snu.ac.kr}\enspace\textperiodcentered\enspace
\link{https://jeongseoko.github.io/}{jeongseoko.github.io}\enspace\textperiodcentered\enspace
\link{https://scholar.google.com/citations?user=WV5svAIAAAAJ}{Google Scholar}\enspace\textperiodcentered\enspace
\link{https://github.com/jeongseokO}{GitHub}\enspace\textperiodcentered\enspace
\link{https://www.linkedin.com/in/\%EC\%A0\%95\%EC\%84\%9D-\%EC\%98\%A4-438528255}{LinkedIn}\par

% ---------------------------------------------------------------- interests
\section{Research Interests}
Efficient and trustworthy large language models: LLM architecture and long-context inference (KV cache, linear and hybrid attention, structured pruning), reliability (self-consistency, hallucination mitigation), and prompting and inference-time strategies.

% ---------------------------------------------------------------- education
\section{Education}
\textbf{Ph.D. in Data Science}, Seoul National University \sub{\textperiodcentered\ Advisor: Prof. Jay-Yoon Lee}\when{Sep 2026 -- Present}\par
\textbf{M.S. in Data Science}, Seoul National University \sub{\textperiodcentered\ Advisor: Prof. Jay-Yoon Lee}\when{Mar 2023 -- Aug 2025}\par
\textbf{B.S. in Industrial Engineering}, Hanyang University\when{Mar 2016 -- Aug 2022}\par

% ---------------------------------------------------------------- publications
\section{Publications}
\textbf{Conference Papers}\hfill\sub{\small * corresponding author}
\begin{pubs}
  \item[{[C1]}] \textbf{Latent Self-Consistency for Reliable Majority-Set Selection in Short- and Long-Answer Reasoning}\\
    \textbf{Jungsuk Oh}, Jay-Yoon Lee*\hfill
    \link{https://doi.org/10.1609/aaai.v40i38.40536}{[Paper]}
    \link{https://arxiv.org/abs/2508.18395}{[arXiv]}
    \link{https://github.com/jeongseokO/LatentSC}{[Code]}
    \link{https://jeongseoko.github.io/research/lsc/}{[Project]}\\
    \textit{AAAI Conference on Artificial Intelligence (AAAI-26)}, 2026. DOI: 10.1609/aaai.v40i38.40536\\
    \sub{Picks the answer most sampled responses agree on by comparing learned summary-token embeddings, so self-consistency also works for long, free-form answers, at most 0.9\% extra inference time.}
\end{pubs}
\textbf{Preprints}
\begin{pubs}
  \item[{[P1]}] \textbf{Shallow Prefill, Deep Decoding: Efficient Long-Context Inference via Layer-Asymmetric KV Visibility}\\
    \textbf{Jungsuk Oh}, Hyeseo Jeon, Hyunjune Ji, Kyongmin Kong, Jay-Yoon Lee*\hfill
    \link{https://arxiv.org/abs/2605.06105}{[arXiv]}
    \link{https://jeongseoko.github.io/research/speed/}{[Project]}\\
    \textit{arXiv:2605.06105}, 2026\\
    \sub{Keeps prompt KV states only in the lower layers while decoding at full depth. On Llama-3.1-8B: 51.2 vs.\ 51.4 average score against the full model, with 33\% faster first token, 22\% faster decoding, and 25\% less KV memory at 128K context.}
\end{pubs}

% ---------------------------------------------------------------- experience
\section{Research Experience}
\textbf{AI Researcher}, Seoul National University R\&DB Foundation (SKI-ML Lab)\when{Sep 2025 -- Aug 2026}
\begin{itemize}
  \item Led the SKI-ML Lab side of the MFDS translation project as head researcher (below).
  \item Researched efficient long-context inference for LLMs, which led to the SPEED preprint [P1].
\end{itemize}
\vspace{4pt}
\textbf{Head Researcher}, MFDS Korean--English Translation Project\when{Feb 2025 -- Nov 2026}\\
\sub{Translation model for drug regulatory documents \textperiodcentered\ SKI-ML Lab $\times$ Ministry of Food and Drug Safety (MFDS)}\\
\sub{PI: Prof. Howard Lee \textperiodcentered\ Co-PI: Prof. Jay-Yoon Lee}
\begin{itemize}
  \item Led model selection and training: compared more than eight candidate models and fine-tuned VaivGem-9B with LoRA SFT and GRPO.
  \item Compared four prompt templates for inference and adopted retrieval few-shot and segment-level prompting.
  \item Built a post-editing loop that checks each translation (length, placeholders, copying, quality estimation) and repairs or flags failures.
  \item Delivered the final model, now hosted on a server inside MFDS.
\end{itemize}

% ---------------------------------------------------------------- teaching
\section{Teaching}
\textbf{Teaching Assistant}, Project Course, Samsung AI Expert Program\when{Jul 2024 -- Aug 2024}
\begin{itemize}
  \item Taught project-based sessions on LLM inference optimization for latency and memory, and guided participants through end-to-end RAG projects (data preprocessing, retriever training, generation).
\end{itemize}

% ---------------------------------------------------------------- service and skills
\section{Service}
\textbf{Reviewer}: ICLR 2026, ICML 2026, AAAI 2026, AAAI 2027

\section{Skills}
\begin{tabularx}{\linewidth}{@{}l X@{}}
  \textbf{ML \& systems} & PyTorch, Hugging Face Transformers, vLLM, PEFT/LoRA, TRL, Slurm\\
  \textbf{Languages} & Python, C/C++, JavaScript, \LaTeX\\
\end{tabularx}

\end{document}
